\documentclass[a4paper,12pt,twoside]{memoir}

% Castellano
\usepackage[spanish,es-tabla]{babel}
\selectlanguage{spanish}
\usepackage[utf8]{inputenc}
\usepackage[T1]{fontenc}
\usepackage{lmodern} % scalable font
\usepackage{microtype}
\usepackage{placeins}

\RequirePackage{booktabs}
\RequirePackage[table]{xcolor}
\RequirePackage{xtab}
\RequirePackage{multirow}

% Links
\PassOptionsToPackage{hyphens}{url}\usepackage[colorlinks]{hyperref}
\hypersetup{
	allcolors = {red}
}

% Ecuaciones
\usepackage{amsmath}

% Rutas de fichero / paquete
\newcommand{\ruta}[1]{{\sffamily #1}}

% Párrafos
\nonzeroparskip

% Huérfanas y viudas
\widowpenalty100000
\clubpenalty100000

% Evitar solapes en el header
\nouppercaseheads

% Imagenes
\usepackage{graphicx}
\newcommand{\imagen}[2]{
	\begin{figure}[!h]
		\centering
		\includegraphics[width=0.9\textwidth]{#1}
		\caption{#2}\label{fig:#1}
	\end{figure}
	\FloatBarrier
}

\newcommand{\imagenflotante}[2]{
	\begin{figure}%[!h]
		\centering
		\includegraphics[width=0.9\textwidth]{#1}
		\caption{#2}\label{fig:#1}
	\end{figure}
}



% El comando \figura nos permite insertar figuras comodamente, y utilizando
% siempre el mismo formato. Los parametros son:
% 1 -> Porcentaje del ancho de página que ocupará la figura (de 0 a 1)
% 2 --> Fichero de la imagen
% 3 --> Texto a pie de imagen
% 4 --> Etiqueta (label) para referencias
% 5 --> Opciones que queramos pasarle al \includegraphics
% 6 --> Opciones de posicionamiento a pasarle a \begin{figure}
\newcommand{\figuraConPosicion}[6]{%
  \setlength{\anchoFloat}{#1\textwidth}%
  \addtolength{\anchoFloat}{-4\fboxsep}%
  \setlength{\anchoFigura}{\anchoFloat}%
  \begin{figure}[#6]
    \begin{center}%
      \Ovalbox{%
        \begin{minipage}{\anchoFloat}%
          \begin{center}%
            \includegraphics[width=\anchoFigura,#5]{#2}%
            \caption{#3}%
            \label{#4}%
          \end{center}%
        \end{minipage}
      }%
    \end{center}%
  \end{figure}%
}

%
% Comando para incluir imágenes en formato apaisado (sin marco).
\newcommand{\figuraApaisadaSinMarco}[5]{%
  \begin{figure}%
    \begin{center}%
    \includegraphics[angle=90,height=#1\textheight,#5]{#2}%
    \caption{#3}%
    \label{#4}%
    \end{center}%
  \end{figure}%
}
% Para las tablas
\newcommand{\otoprule}{\midrule [\heavyrulewidth]}
%
% Nuevo comando para tablas pequeñas (menos de una página).
\newcommand{\tablaSmall}[5]{%
 \begin{table}
  \begin{center}
   \rowcolors {2}{gray!35}{}
   \begin{tabular}{#2}
    \toprule
    #4
    \otoprule
    #5
    \bottomrule
   \end{tabular}
   \caption{#1}
   \label{tabla:#3}
  \end{center}
 \end{table}
}

%
%Para el float H de tablaSmallSinColores
\usepackage{float}

%
% Nuevo comando para tablas pequeñas (menos de una página).
\newcommand{\tablaSmallSinColores}[5]{%
 \begin{table}[H]
  \begin{center}
   \begin{tabular}{#2}
    \toprule
    #4
    \otoprule
    #5
    \bottomrule
   \end{tabular}
   \caption{#1}
   \label{tabla:#3}
  \end{center}
 \end{table}
}

\newcommand{\tablaApaisadaSmall}[5]{%
\begin{landscape}
  \begin{table}
   \begin{center}
    \rowcolors {2}{gray!35}{}
    \begin{tabular}{#2}
     \toprule
     #4
     \otoprule
     #5
     \bottomrule
    \end{tabular}
    \caption{#1}
    \label{tabla:#3}
   \end{center}
  \end{table}
\end{landscape}
}

%
% Nuevo comando para tablas grandes con cabecera y filas alternas coloreadas en gris.
\newcommand{\tabla}[6]{%
  \begin{center}
    \tablefirsthead{
      \toprule
      #5
      \otoprule
    }
    \tablehead{
      \multicolumn{#3}{l}{\small\sl continúa desde la página anterior}\\
      \toprule
      #5
      \otoprule
    }
    \tabletail{
      \hline
      \multicolumn{#3}{r}{\small\sl continúa en la página siguiente}\\
    }
    \tablelasttail{
      \hline
    }
    \bottomcaption{#1}
    \rowcolors {2}{gray!35}{}
    \begin{xtabular}{#2}
      #6
      \bottomrule
    \end{xtabular}
    \label{tabla:#4}
  \end{center}
}

%
% Nuevo comando para tablas grandes con cabecera.
\newcommand{\tablaSinColores}[6]{%
  \begin{center}
    \tablefirsthead{
      \toprule
      #5
      \otoprule
    }
    \tablehead{
      \multicolumn{#3}{l}{\small\sl continúa desde la página anterior}\\
      \toprule
      #5
      \otoprule
    }
    \tabletail{
      \hline
      \multicolumn{#3}{r}{\small\sl continúa en la página siguiente}\\
    }
    \tablelasttail{
      \hline
    }
    \bottomcaption{#1}
    \begin{xtabular}{#2}
      #6
      \bottomrule
    \end{xtabular}
    \label{tabla:#4}
  \end{center}
}

%
% Nuevo comando para tablas grandes sin cabecera.
\newcommand{\tablaSinCabecera}[5]{%
  \begin{center}
    \tablefirsthead{
      \toprule
    }
    \tablehead{
      \multicolumn{#3}{l}{\small\sl continúa desde la página anterior}\\
      \hline
    }
    \tabletail{
      \hline
      \multicolumn{#3}{r}{\small\sl continúa en la página siguiente}\\
    }
    \tablelasttail{
      \hline
    }
    \bottomcaption{#1}
  \begin{xtabular}{#2}
    #5
   \bottomrule
  \end{xtabular}
  \label{tabla:#4}
  \end{center}
}



\definecolor{cgoLight}{HTML}{EEEEEE}
\definecolor{cgoExtralight}{HTML}{FFFFFF}

%
% Nuevo comando para tablas grandes sin cabecera.
\newcommand{\tablaSinCabeceraConBandas}[5]{%
  \begin{center}
    \tablefirsthead{
      \toprule
    }
    \tablehead{
      \multicolumn{#3}{l}{\small\sl continúa desde la página anterior}\\
      \hline
    }
    \tabletail{
      \hline
      \multicolumn{#3}{r}{\small\sl continúa en la página siguiente}\\
    }
    \tablelasttail{
      \hline
    }
    \bottomcaption{#1}
    \rowcolors[]{1}{cgoExtralight}{cgoLight}

  \begin{xtabular}{#2}
    #5
   \bottomrule
  \end{xtabular}
  \label{tabla:#4}
  \end{center}
}




\graphicspath{ {./img/} }

% Capítulos
\chapterstyle{bianchi}
\newcommand{\capitulo}[2]{
	\setcounter{chapter}{#1}
	\setcounter{section}{0}
	\setcounter{figure}{0}
	\setcounter{table}{0}
	\chapter*{#2}
	\addcontentsline{toc}{chapter}{#2}
	\markboth{#2}{#2}
}

% Apéndices
\renewcommand{\appendixname}{Apéndice}
\renewcommand*\cftappendixname{\appendixname}

\newcommand{\apendice}[1]{
	%\renewcommand{\thechapter}{A}
	\chapter{#1}
}

\renewcommand*\cftappendixname{\appendixname\ }

% Formato de portada
\makeatletter
\usepackage{xcolor}
\newcommand{\tutor}[1]{\def\@tutor{#1}}
\newcommand{\course}[1]{\def\@course{#1}}
\definecolor{cpardoBox}{HTML}{E6E6FF}
\def\maketitle{
  \null
  \thispagestyle{empty}
  % Cabecera ----------------
\noindent\includegraphics[width=\textwidth]{cabecera}\vspace{1cm}%
  \vfill
  % Título proyecto y escudo informática ----------------
  \colorbox{cpardoBox}{%
    \begin{minipage}{.8\textwidth}
      \vspace{.5cm}\Large
      \begin{center}
      \textbf{TFG del Grado en Ingeniería Informática}\vspace{.6cm}\\
      \textbf{\LARGE\@title{}}
      \end{center}
      \vspace{.2cm}
    \end{minipage}

  }%
  \hfill\begin{minipage}{.20\textwidth}
    \includegraphics[width=\textwidth]{escudoInfor}
  \end{minipage}
  \vfill
  % Datos de alumno, curso y tutores ------------------
  \begin{center}%
  {%
    \noindent\LARGE
    Presentado por \@author{}\\ 
    en Universidad de Burgos --- \@date{}\\
    Tutor: \@tutor{}\\
    Co-tutor: Sandra Rodríguez Arribas\\
    Tutor de Empresa: Julia Zuara Jiménez\\
  }%
  \end{center}%
  \null
  \cleardoublepage
  }
\makeatother


% Datos de portada
\title{Second Hand Chain \\Documentación Técnica}
\author{Áxel Rubio González}
\tutor{Jesús Manuel Maudes Raedo}
\date{\today}

\begin{document}

\maketitle



\cleardoublepage



%%%%%%%%%%%%%%%%%%%%%%%%%%%%%%%%%%%%%%%%%%%%%%%%%%%%%%%%%%%%%%%%%%%%%%%%%%%%%%%%%%%%%%%%



\frontmatter


\clearpage

% Indices
\tableofcontents

\clearpage

\listoffigures

\clearpage

\listoftables

\clearpage

\mainmatter

\appendix

\apendice{Plan de Proyecto Software}

\section{Introducción}

En este capítulo se va a exponer cómo se ha planeado el proyecto y cómo se ha llevado a cabo dentro de cada uno de los apartados.

\section{Planificación temporal}
La gestión de este proyecto se ha realizado empleando la metodología \textit{Scrum}. Por ello, se celebraba con los tutores el \textit{sprint review}, una reunión para revisar todo el trabajo realizado y resolver posibles dudas, esta era realizada con una frecuencia de dos semanas, coincidiendo con el final e inicio de \textit{sprint} .

Para la clasificación de \textit{issues} se ha empleado un sistema que los divide por las siguientes etiquetas para poder facilitar mejor su identificación:
\begin{enumerate}
    \item \textit{Backend}: Empleada para \textit{issues} relacionadas con el \textit{backend}, predomina durante la primera mitad del proyecto, fase en la que se estaban creando los contratos digitales.
    \item \textit{Blockchain}: Se ha usado para \textit{issues} relacionadas con la tecnología \textit{Blockchain}, predomina durante la primera mitad del proyecto, suele ir asociada a la de \textit{Backend}.
    \item \textit{Bug}: Al final del proyecto se pudieron detectar varios pequeños errores en el \textit{Frontend}. Para dar más visibilidad a esas tareas se creó esta etiqueta.
    \item \textit{Frontend}: Durante la segunda mitad del proyecto el desarrollo se centra en el \textit{Frontend} por lo que es cuando comienza a tomar protagonismo.
    \item \textit{Management}: Etiqueta relacionada con actividades de la gestión de tareas.
    \item \textit{Research}: Para tareas que incluyen investigación para conocer cómo realizar dicha tarea.
    \item \textit{Software engineering}: Para las actividades realizadas durante las fases de diseño inicial de la aplicación.
\end{enumerate}

\subsection{Sprint 1: Inicio del proyecto}

\imagen{anexos/sprint1}{Tareas finalizadas en el Sprint 1}

En este primer sprint transcurrió entre el 24-02-2023 y el 10-03-2023  se realizo una primera toma de contacto con el proyecto, consistió en investigar el funcionamiento de \textit{Blockchain}. Además, se comenzó a aprender a usar \textit{GitLab} para realizar la gestión de tareas durante el proyecto. Finalmente, se realizaron unos primeros diagramas UML a partir de los requisitos ideados en la reunión incial del proyecto. Podemos ver todas estas tareas con sus etiquetas en la figura \ref{fig:anexos/sprint1}.

\subsection{Sprint 2: Configuraciones iniciales y primera iteración del \textit{smart contract}}

\imagen{anexos/sprint2}{Tareas finalizadas en el Sprint 2}

En este segundo \textit{sprint}, el cual transcurrió entre el 11-03-2023 y el 24-03-2023, se comenzó a configurar el repositorio de GitLab, al cual solo se podía acceder mediante clave SSH, por ello tuve que configurar una y después resolver unos cuantos problemas que me daban las cuentas de \textit{git}.

A partir de entonces, ya con el repositorio configurado, comencé a formarme en \textit{Solidity}, un lenguaje que debido a las características de \textit{Blockchain} tiene una serie de particularidades que deben ser bien comprendidas para empezar a emplearlo. Después de realizar esta tarea, se creó un primer \textit{smart contract} que permitía el registro de teléfonos y definía inicialmente cómo se iban a almacenar estos datos.

\subsection{Sprint 3: Adaptación del \textit{smart contract} al estándar \textit{ERC721}}

Transcurrido entre el 25-03-2023 y el 14-04-2023 se continuó con el desarrollo del \textit{smart contract}, se implementaron nuevos \textit{getters}, proceso que es distinto al empleado en otros lenguajes ya que se debe guardar en en local el \textit{array} para estos métodos \textit{view} (sin coste al solo consultar información) y por ello debes almacenar el número de objetos que vas a añadir a esa lista que debe conocer de antemano el tamaño que va a tener al instanciarla, no es dinámica como los \textit{arrays} en \textit{storage}.

\imagen{anexos/sprint3}{Tareas finalizadas en el Sprint 3}

Además, aunque de esto no se pudo hacer \textit{commit} al detectar un error que se tuvo que preguntar a los tutores, se comenzó con la implementación de la interfaz \textit{IERC721}, el estándar para gestionar traspasos de \textit{NFTs}.

\subsection{Sprint 4: Continuación de la adaptación al estándar \textit{ERC721}}

Entre el 15-04-2023 y el 21-04-2023, este sprint fue solo de una semana (el anterior de tres) por problemas de horarios por la semana santa, se resolvió el error descrito anteriormente que se debía a un problema con las versiones del compilador de \textit{Solidity}.

\imagen{anexos/sprint4}{Tareas finalizadas en el Sprint 4}

Además, se finalizó del desarrollo de las características básicas necesarias a implementar en los \textit{smart contract}. Para la implementación del estándar \textit{ERC721} se crearon otros métodos que no se iban a usar en Second Hand Chain, principalmente los correspondientes a dar permisos a terceros para transferir el token, para dejar la puerta abierta en un futuro a su uso por otras aplicaciones.

Hasta esta fase, se había usado la consola de \textit{Truffle} y \textit{Ganache} para comprobar el correcto funcionamiento de estos \textit{smart contract}, es decir, funcionaban correctamente n el entorno de desarrollo local.

\subsection{Sprint 5: Inicio del desarrollo del \textit{frontend}}


En este sprint, entre el 22-04-2023 y el 5-05-2023,  se desarrolla todo lo relacionado con el despliegue del \textit{backend} en \textit{Sepolia} y el inicio del \textit{frontend}.

\imagen{anexos/sprint5}{Tareas finalizadas en el Sprint 5}

Se toma la decisión de migrar en este momento el contrato a la \textit{testnet} para facilitar el desarrollo y evitar posibles problemas de integración si el despliegue se produce con el \textit{frontend} ya terminado.

En consecuencia, se inicia el proyecto de \textit{React}, en el cual se empieza implementando un \textit{custom hook} para obtener los móviles y se empieza con el diseño de un modal para crear teléfonos.

\subsection{Sprint 6: Continuamos con el \textit{frontend}}



En el sprint que transcurre entre el 6-05-2023 y el 26-05-2023 se continúa con el desarrollo del \textit{frontend},  tras el cual se va a permitir que los usuarios puedan consultar los teléfonos a la venta y los teléfonos propios, siempre y cuando se tenga \textit{wallet} y teléfonos registrados.

\imagen{anexos/sprint6}{Tareas finalizadas en el Sprint 6}

El otro tipo de usuario, el que no dispone de \textit{wallet}, podrá consultar los teléfonos a la venta aunque no podrá ni crear ni comprar.

\subsection{Sprint 7: Llegamos al producto mínimo viable}



En este sprint se alcanzó el producto mínimo viable, transcurrió ente el 27-05-2023 y el 9-06-2023. Se implementaron todas las funciones necesarias para cumplir todos los requisitos funcionales de la aplicación. 
Se implementó la compra venta de teléfonos, la posibilidad de actualizar el precio de los productos que están ya a la venta y retirar del mercado. Además se implementó un nuevo modal que permite a los usuarios controlar todas estas funciones anteriores por teléfono y consultar el historial de precios y ventas previas.

\imagen{anexos/sprint7}{Tareas finalizadas en el Sprint 7}

En consecuencia, se produce el primer despliegue de la aplicación en \textit{Vercel} para comprobar su correcto funcionamiento una vez ya desplegada.

\subsection{Sprint 8: Grandes Mejoras a nivel de \textit{UX} y nuevas características}

En este sprint final transcurrido entre el 10-06-2023 y el 23-06-2023, se decide añadir a la aplicación la posibilidad de registrar imágenes de los teléfonos, ese añadido se realiza respetando las ventajas de \textit{Blockchain}. Por ello este sprint está focalizado principalmente en esta característica que requiere de la adaptación del contrato para poder registrar el CID y en la adaptación del \textit{frontend} para esta característica.

\imagen{anexos/sprint81}{Tareas finalizadas en el Sprint 8 (Parte 1)}
\imagen{anexos/sprint82}{Tareas finalizadas en el Sprint 8 (Parte 2)}

Se mejora el funcionamiento del botón de la parte superior derecha que indica el estado de conexión de la aplicación con la \textit{wallet}, haciéndolo mucho más comprensible para el usuario, de esta forma ahora el usuario conoce en cualquier momento si tiene \textit{Metamask} instalado, si lo tiene pero debe iniciar sesión o si ya está conectado.

Surgieron una serie de \textit{bugs}, principalmente estéticos por actualizar \textit{DaisyUi} a una nueva versión que había salido en Junio que tuvieron que ser resueltos y otros fallos menores.

Finalmente, se realizaron mejoras para garantizar la seguridad de la aplicación, se creó un fichero \textit{.env} para garantizar la seguridad de las distintas \textit{API KEYS}, además estas claves fueron regeneradas. La aplicación se volvió a desplegar y se creo un repositorio paralelo al original en \textit{GitHub} para poder hacer público el proyecto.

\section{Estudio de viabilidad}

En esta sección se va a estudiar la viabilidad tanto económica como legal del proyecto, para ello, se va a hacer un repaso de las tecnologías empleadas para ver que costes conlleva y las licencias empleadas.

\subsection{Viabilidad económica}
En este apartado se estudiarán los costes económicos del proyecto y posibles formas de monetizar el desarrollo.

\subsubsection{Costes del proyecto}

En primer lugar vamos a considerar los costes del proyecto relacionados con el personal:

\begin{itemize}
    \item En total se han empleado alrededor de 250 horas para realizar el proyecto, entre realización de las tareas en sí e investigación, ello da una media de 14 horas a la semana.
    \item Las principales tecnologías empleadas han sido \textit{Solidity} y \textit{React}, sus sueldos son respectivamente de 70.368 USD\footnote{United States Dolar} y 69.480 USD, para estimar un sueldo para nuestro desarrollador vamos a calcular la media de ambas para estimar un sueldo medio que quiera aceptar nuestro desarrollador, serán 69.924\$. Estos sueldos han sido extraídos de la encuesta anual de \textit{Stack Overflow}\cite{stackoverflowsurvey}. El sueldo bruto anual en euros será de 64.053 euros a la tasa de cambio actual por la que 1 USD son 0,92 euros. Esto nos daría como resultado unos 35 euros por hora brutos a 1820 horas de trabajo anuales (40 semanales).

 $$ 14 \frac{horas}{ 1 semana}\times35 \frac{euros}{ 1 hora}\times4\frac{semanas}{ 1 mes}= 1960 euros\hspace{0.5em}al\hspace{0.5em}mes  $$

\item A continuación, debemos pagar los impuestos para calcular el coste total del trabajador, datos extraídos de \cite{impuestosSS}:
\begin{itemize}
    \item FOGASA: 0,2\%
    \item Desempleo: 5,5\%
    \item Formación: 0,6\%
    \item Contingencias: 23,6\%
\end{itemize}

$$\frac{1960\frac{euros}{ 1 mes}}{1-(0.002+0.055+0.006+0.236)}=2796\text{\euro}\hspace{0.5em}al\hspace{0.5em}mes$$
\end{itemize}
Además, se deben añadir los costes del equipo informático y su software:
\begin{itemize}
    \item  Un ordenador de sobremesa con un coste total de unos 2200 euros después diversas mejoras, aunque el original es de hace 7 años, actualmente tiene las siguientes características:
    \begin{itemize}
        \item Procesador: AMD Ryzen 5 5600X 6 núcleos.
        \item Gráfica: Nvidia GTX 1070.
        \item RAM: 16GB DDR4.
        \item SSD: 750GB.
        \item HD: 2Tb.
    \end{itemize}
$$\frac{2200\hspace{0.5em}euros}{7\hspace{0.5em}años}=314\hspace{0.5em}euros\hspace{0.5em}al\hspace{0.5em}año$$

$$314euros\frac{4\hspace{0.5em} meses}{12\hspace{0.5em} meses}=104euros\hspace{0.5em}en\hspace{0.5em} total\hspace{0.5em} por\hspace{0.5em} hardware$$
\item  De software solo se ha pagado Windows 10, hace 7 años también:
    $$\frac{130\hspace{0.5em}euros}{7\hspace{0.5em}años}=19\hspace{0.5em}euros\hspace{0.5em}al\hspace{0.5em}año$$
    $$19euros\frac{4\hspace{0.5em} meses}{12\hspace{0.5em} meses}=6.2euros\hspace{0.5em}en\hspace{0.5em} total\hspace{0.5em} por\hspace{0.5em} software$$
\end{itemize}


En consecuencia, podemos concluir que los costes totales del proyecto en estos cuatro meses han sido de:
    \begin{itemize}
        \item Personal: 11184 euros.
        \item Hardware: 104 euros.
        \item Software: 6,2 euros.
        \item Otros (electricidad y conexión a Internet): 130 euros.
    \end{itemize}

En total suma unos 11424 euros.

\subsubsection{Costes del proyecto}
Actualmente este proyecto al no haber sido desplegado en la red principal y además al ser distribuido con licencia MIT no va a recaudar dinero. No obstante, si lo deseáramos, dispondríamos de las siguientes alternativas:
    \begin{itemize}
        \item Cobrar una tasa por transacción a nivel de contratos. El contrato enviaría esta tasa a la \textit{wallet} definida. Además, los usuarios podrían consultar el funcionamiento esta tasa a nivel de código.
        \item De forma similar a la anterior, cobrar una tasa por usar el \textit{frontend}.
        \item Proponer el proyecto como \textit{open source} y financiarlo mediante una campaña de \textit{crowdfunding}.
    \end{itemize}

\subsection{Viabilidad legal}

\tablaSmall{Licencias empleadas en el proyecto}
{l r r r}{licenses}
{\textbf{Herramienta} & \textbf{Licencia}  \\}{
  React & MIT  \\
  Tailwinds css & MIT  \\
  Daisy UI & MIT  \\
  Solidity & GNU  \\
  GIT & GNU \\
  NFT.Storage & MIT y Apache License 2.0 \\
  Truffle Suite & MIT \\
  Web3.js & GNU \\
  Vite & MIT \\
  TypeScript & Apache License 2.0 \\
  Metamask & Non-Commercial Use* \\
  Node.js & MIT \\
  Ethereum & GNU \\
}

Cómo se puede ver en la tabla \ref{tabla:licenses} todas las licencias permiten su uso para proyectos del tipo de Second Hand Chain de forma completamente libre.

La única particularidad sería la de \textit{Metamask}, con cuya compañía se necesitaría llegar a un acuerdo para emplear su software en caso de que se usase para fines comerciales o se superase los 10000 usuarios mensuales, cifras que no vamos a alcanzar en ningún caso.





\apendice{Especificación de Requisitos}

\section{Introducción}

En este apéndice se van a listar los requisitos que satisface la aplicación que forma parte de este proyecto. Aunque los requisitos se definieron en la fase de diseño de la aplicación este listado ha sido ampliado según se desarrollaba la aplicación.


\section{Objetivos generales}
A continuación se adjunta un listado de los objetivos principales de la aplicación:
\begin{itemize}
    \item Crear un \textit{marketplace} que permita a los usuarios vender y comprar \textit{NFT}.
    \item Desplegar los contratos en la \textit{Blockchain} \textit{Ethereum} o una de sus \textit{testnets}.
    \item Aportar a los usuarios información que por sus características de inmutabilidad respaldada por \textit{Ethereum} aporte valor a sus productos.
    \item Acercar las tecnologías \textit{Blockchain} a un mayor numero de usuarios.
    \item Crear un aplicación web que aporte una experiencia lo más accesible posible a estas tecnologías \textit{Blockchain}.
    
\end{itemize}

\section{Catalogo de requisitos}

En esta sección se incluye un listado con todos los requisitos que debe cumplir la aplicación:

\subsection{Requisitos Funcionales}
\begin{itemize}
    \item \textbf{RF1-Gestión de \textit{wallet}}: El usuario deberá poder iniciar sesión en la aplicación desde la extensión de \textit{Metamask} de su navegador, si tiene la \textit{wallet} instalada le aparecerá directamente para iniciar sesión, si no, se mostrarán los correspondientes mensajes de aviso.
    \item \textbf{RF2-Gestión de usuario sin \textit{wallet}}: Un usuario que no cuente con \textit{wallet} instalada podrá usar la aplicación para consultar teléfonos a la venta y características. Además de la \textit{landing page} en la cual podrá consultar cómo añadir esta extensión si lo deseara.
    \item \textbf{RF3-Creación de teléfonos}: Se debe poder registrar los teléfonos en \textit{Blockchain}, para ello el usuario deberá introducir todos sus datos además de una imagen. No se podrá registrar un teléfono con un IMEI ya registrado.
    \item \textbf{RF4-Gestión de teléfonos}: Los usuarios registrados deben poder ver los teléfonos que les pertenecen, además, podrán acceder al mercado de teléfonos en venta.
    \item \textbf{RF5-Gestión de ventas}: Los usuarios deben ser capaces de gestionar el estado en que se encuentren sus teléfonos, si se venden o no.
    \begin{itemize}
        \item \textbf{RF5.1-Puesta a la venta}: Los usuarios pueden poner a la venta los teléfonos previamente registrados.
        \item \textbf{RF5.2-Cambio de precio}: Dado un teléfono que está en venta su propietario podrá cambiar su precio de venta.
        \item \textbf{RF5.2-Retirada de la venta}: Dado un teléfono que está en venta su propietario podrá retirarlo de la venta.
    \end{itemize}
    \item \textbf{RF6-\textbf{Seguimiento de transacciones por \textit{Etherscan}}}: Se requiere que por cada por cada operación que implique realizar una transacción en \textit{Blockchain}, el usuario, tras firmar dicha transacción, dispondrá de un enlace que le permitirá seguir su estado a través de \textit{Etherscan}.
    \item \textbf{RF7-Compra de teléfonos}: Los usuarios deben poder comprar los teléfonos, para realizar esta operación aportarán el \textit{Ether} correspondiente a la transacción , siempre y cuando este teléfono no les pertenezca ya.
\end{itemize}

\subsection{Requisitos No Funcionales}
\begin{itemize}
    \item \textbf{RNF1-Seguridad}: Se requiere que la aplicación ofrezca un entorno seguro para el usuario y no suponga ningún riesgo para sus datos el uso de ella.
    \item \textbf{RNF2-Disponibilidad}: Se requiere que la aplicación esté disponible constantemente excepto en sus periodos de mantenimiento.
    \item \textbf{RNF3-Usabilidad}: Se requiere que la aplicación ofrezca la mejor experiencia posible al usuario. Además, es necesaria la correcta visualización de la aplicación desde teléfonos móviles, por lo tanto, los elementos en pantalla que sean necesarios se adaptarán dinámicamente para estas pantallas.
    \item \textbf{RNF4-Confiabilidad}: Se requiere que la aplicación cumpla todos los requisitos que se establecen en estos anexos.
    \item \textbf{RNF5-Eficiencia}: Se requiere que la aplicación haga un uso eficiente de los recursos de la máquina que lo vaya a ejecutar.
    \item \textbf{RNF6-Privacidad}: Se requiere que la aplicación haga un uso responsable de los datos de los usuarios, en este caso de su clave privada, dato que no debe ser conocido por nadie excepto el propio usuario.
    \item \textbf{RNF7-Mantenibilidad}: La aplicación debe aplicar las mejores prácticas posibles para facilitar el mantenimiento durante todo el ciclo de vida de la aplicación.
\end{itemize}


\section{Especificación de requisitos}

\subsection{Diagrama de casos de uso}
\imagen{anexos/dicu}{Diagrama de casos de uso}

% Caso de Uso 1 -> Consultar Experimentos.
\begin{table}[p]
	\centering
	\begin{tabularx}{\linewidth}{ p{0.21\columnwidth} p{0.71\columnwidth} }
		\toprule
		\textbf{CU-1}    & \textbf{Inicio de sesión en \textit{Metamask}}\\
		\toprule
		\textbf{Versión}              & 1.0    \\
		\textbf{Autor}                & Áxel Rubio González \\
		\textbf{Requisitos asociados} & RF-1 \\
		\textbf{Descripción}          & En este caso de uso el usuario iniciará sesión en \textit{Metamask} \\
		\textbf{Precondición}         & El usuario no ha iniciado sesión en \textit{Metamask} y tiene la extensión instalada \\
		\textbf{Acciones}             &
		\begin{enumerate}
			\def\labelenumi{\arabic{enumi}.}
			\tightlist
			\item El usuario accede a la web de Second Hand Chain.
			\item Le aparece una ventana emergente para iniciar sesión en \textit{Metamask}.
                \item El usuario introduce su contraseña.
                \item El usuario selecciona la wallet con la que quiere operar.
		\end{enumerate}\\
		\textbf{Postcondición}        & Que \textit{Metamask} mantenga esa sesión iniciada \\
		\textbf{Excepciones}          & \textit{Metamask} no se encuentra instalado o se ha cerrado la ventana emergente \\
		\textbf{Importancia}          & Alta \\
		\bottomrule
	\end{tabularx}
	\caption{CU-1 Inicio de sesión en \textit{Metamask}.}
\end{table}

\begin{table}[p]
	\centering
	\begin{tabularx}{\linewidth}{ p{0.21\columnwidth} p{0.71\columnwidth} }
		\toprule
		\textbf{CU-2}    & \textbf{Usuario sin \textit{Metamask} accede a los teléfonos en venta}\\
		\toprule
		\textbf{Versión}              & 1.0    \\
		\textbf{Autor}                & Áxel Rubio González \\
		\textbf{Requisitos asociados} & RF2 \\
		\textbf{Descripción}          & En este caso de uso el usuario sin \textit{Metamask} consultará los teléfonos a la venta.  \\
		\textbf{Precondición}         & El usuario no ha iniciado sesión en \textit{Metamask} o no tiene la extensión instalada. \\
		\textbf{Acciones}             &
		\begin{enumerate}
			\def\labelenumi{\arabic{enumi}.}
			\tightlist
			\item El usuario accede a la web de Second Hand Chain.
			\item Hace clic en teléfonos en venta.
                \item Si está en pantalla pequeña clic en el \textit{dropdown} y después en en teléfonos en venta.
                \item El usuario puede visualizar los teléfonos a la venta.
                \item El usuario accede a un teléfono y no podrá comprarlo, le aparece un mensaje diciendo que inicie sesión justamente en el botón.
		\end{enumerate}\\
		\textbf{Postcondición}        & Ninguna \\
		\textbf{Excepciones}          & Ninguna \\
		\textbf{Importancia}          & Alta \\
		\bottomrule
	\end{tabularx}
	\caption{CU-2 Usuario sin \textit{Metamask} accede a los teléfonos en venta.}
\end{table}

\begin{table}[p]
	\centering
	\begin{tabularx}{\linewidth}{ p{0.21\columnwidth} p{0.71\columnwidth} }
		\toprule
		\textbf{CU-3}    & \textbf{Creación de teléfonos}\\
		\toprule
		\textbf{Versión}              & 1.0    \\
		\textbf{Autor}                & Áxel Rubio González \\
		\textbf{Requisitos asociados} & RF3 y RF6 \\
		\textbf{Descripción}          & En este caso de uso el usuario con \textit{Metamask} registrará un teléfono  \\
		\textbf{Precondición}         & El usuario ha iniciado sesión en \textit{Metamask} \\
		\textbf{Acciones}             &
		\begin{enumerate}
			\def\labelenumi{\arabic{enumi}.}
			\tightlist
			\item El usuario accede a la web de Second Hand Chain.
			\item Hace clic en crear teléfono.
                \item Rellena los datos necesarios y selecciona un archivo.
                \item Hace clic en crear teléfono.
                \item La extensión de \textit{Metamask} le solicita firmar la transacción, lo cual acepta.
                \item Le aparece un mensaje para seguir la transacción en \textit{Etherscan}.
		\end{enumerate}\\
		\textbf{Postcondición}        & El teléfono se registra correctamente en \textit{Blockchain} \\
		\textbf{Excepciones}          & No se ha iniciado sesión, no hay fondos o no se han rellenado los datos \\
		\textbf{Importancia}          & Alta \\
		\bottomrule
	\end{tabularx}
	\caption{CU-3 Creación de teléfonos.}
\end{table}

\begin{table}[p]
	\centering
	\begin{tabularx}{\linewidth}{ p{0.21\columnwidth} p{0.71\columnwidth} }
		\toprule
		\textbf{CU-4}    & \textbf{Gestión de teléfonos}\\
		\toprule
		\textbf{Versión}              & 1.0    \\
		\textbf{Autor}                & Áxel Rubio González \\
		\textbf{Requisitos asociados} & RF4 \\
		\textbf{Descripción}          & En este caso de uso el usuario con \textit{Metamask} consultará los datos de todos los teléfonos  \\
		\textbf{Precondición}         & El usuario ha iniciado sesión en \textit{Metamask} \\
		\textbf{Acciones}             &
		\begin{enumerate}
			\def\labelenumi{\arabic{enumi}.}
			\tightlist
			\item El usuario accede a la web de Second Hand Chain.
			\item Hace clic en sus teléfonos.
                \item Puede consultar los datos del teléfono.
                \item Hace clic en teléfonos en venta.
                \item Puede acceder a todos los teléfonos que se están vendiendo.
                \item Consulta las características de cualquiera de ellos.
		\end{enumerate}\\
		\textbf{Postcondición}        & Ninguna \\
		\textbf{Excepciones}          & Ninguna \\
		\textbf{Importancia}          & Alta \\
		\bottomrule
	\end{tabularx}
	\caption{CU-4 Gestión de teléfonos.}
\end{table}


\begin{table}[p]
	\centering
	\begin{tabularx}{\linewidth}{ p{0.21\columnwidth} p{0.71\columnwidth} }
		\toprule
		\textbf{CU-5}    & \textbf{Puesta a la venta}\\
		\toprule
		\textbf{Versión}              & 1.0    \\
		\textbf{Autor}                & Áxel Rubio González \\
		\textbf{Requisitos asociados} & RF51 y RF6 \\
		\textbf{Descripción}          & En este caso de uso el usuario con \textit{Metamask} podrá poner un teléfono a la venta  \\
		\textbf{Precondición}         & El usuario ha iniciado sesión en \textit{Metamask} y tiene teléfonos \\
		\textbf{Acciones}             &
		\begin{enumerate}
			\def\labelenumi{\arabic{enumi}.}
			\tightlist
			\item El usuario accede a la web de Second Hand Chain.
			\item Hace clic en sus teléfonos.
                \item Selecciona uno.
                \item Define un precio de venta y hace clic en el botón.
                \item La extensión de \textit{Metamask} le solicita firmar la transacción, lo cual acepta.
                \item Le aparece un mensaje para seguir la transacción en \textit{Etherscan}.
		\end{enumerate}\\
		\textbf{Postcondición}        & La puesta a la venta se registra correctamente en \textit{Blockchain} \\
		\textbf{Excepciones}          & No se ha iniciado sesión, no hay fondos o no se han rellenado los datos \\
		\textbf{Importancia}          & Alta \\
		\bottomrule
	\end{tabularx}
	\caption{CU-5 Puesta a la venta de teléfono.}
\end{table}

\begin{table}[p]
	\centering
	\begin{tabularx}{\linewidth}{ p{0.21\columnwidth} p{0.71\columnwidth} }
		\toprule
		\textbf{CU-6}    & \textbf{Cambio de Precio}\\
		\toprule
		\textbf{Versión}              & 1.0    \\
		\textbf{Autor}                & Áxel Rubio González \\
		\textbf{Requisitos asociados} & RF52 y RF6 \\
		\textbf{Descripción}          & En este caso de uso el usuario con \textit{Metamask} podrá poner un teléfono a la venta  \\
		\textbf{Precondición}         & El usuario ha iniciado sesión en \textit{Metamask} y tiene teléfonos en venta \\
		\textbf{Acciones}             &
		\begin{enumerate}
			\def\labelenumi{\arabic{enumi}.}
			\tightlist
			\item El usuario accede a la web de Second Hand Chain.
			\item Hace clic en sus teléfonos.
                \item Selecciona uno en venta.
                \item Define un precio de venta y hace clic en el botón.
                \item La extensión de \textit{Metamask} le solicita firmar la transacción, lo cual acepta.
                \item Le aparece un mensaje para seguir la transacción en \textit{Etherscan}.
		\end{enumerate}\\
		\textbf{Postcondición}        & El cambio de precio se registra correctamente en \textit{Blockchain} \\
		\textbf{Excepciones}          & No se ha iniciado sesión, no hay fondos o no se han rellenado los datos \\
		\textbf{Importancia}          & Alta \\
		\bottomrule
	\end{tabularx}
	\caption{CU-6 Cambio de Precio.}
\end{table}

\begin{table}[p]
	\centering
	\begin{tabularx}{\linewidth}{ p{0.21\columnwidth} p{0.71\columnwidth} }
		\toprule
		\textbf{CU-7}    & \textbf{Retirar de la venta}\\
		\toprule
		\textbf{Versión}              & 1.0    \\
		\textbf{Autor}                & Áxel Rubio González \\
		\textbf{Requisitos asociados} & RF53 y RF6 \\
		\textbf{Descripción}          & En este caso de uso el usuario con \textit{Metamask} podrá poner un teléfono a la venta  \\
		\textbf{Precondición}         & El usuario ha iniciado sesión en \textit{Metamask} y tiene teléfonos en venta \\
		\textbf{Acciones}             &
		\begin{enumerate}
			\def\labelenumi{\arabic{enumi}.}
			\tightlist
			\item El usuario accede a la web de Second Hand Chain.
			\item Hace clic en sus teléfonos.
                \item Selecciona uno en venta.
                \item Se hace click en retirar.
                \item La extensión de \textit{Metamask} le solicita firmar la transacción, lo cual acepta.
                \item Le aparece un mensaje para seguir la transacción en \textit{Etherscan}.
		\end{enumerate}\\
		\textbf{Postcondición}        & La retirada de la venta se registra correctamente en \textit{Blockchain} \\
		\textbf{Excepciones}          & No se ha iniciado sesión, no hay fondos o no se han rellenado los datos \\
		\textbf{Importancia}          & Alta \\
		\bottomrule
	\end{tabularx}
	\caption{CU-7 Retirar de la venta.}
\end{table}

\begin{table}[p]
	\centering
	\begin{tabularx}{\linewidth}{ p{0.21\columnwidth} p{0.71\columnwidth} }
		\toprule
		\textbf{CU-8}    & \textbf{Compra}\\
		\toprule
		\textbf{Versión}              & 1.0    \\
		\textbf{Autor}                & Áxel Rubio González \\
		\textbf{Requisitos asociados} & RF7 y RF6 \\
		\textbf{Descripción}          & En este caso de uso el usuario con \textit{Metamask} comprará un teléfono  \\
		\textbf{Precondición}         & El usuario ha iniciado sesión en \textit{Metamask} y tiene fondos \\
		\textbf{Acciones}             &
		\begin{enumerate}
			\def\labelenumi{\arabic{enumi}.}
			\tightlist
			\item El usuario accede a la web de Second Hand Chain.
			\item Hace clic en teléfonos en venta.
                \item Selecciona uno en venta que no le pertenezca.
                \item Se hace click en comprar.
                \item La extensión de \textit{Metamask} le solicita firmar la transacción, lo cual acepta.
                \item Le aparece un mensaje para seguir la transacción en \textit{Etherscan}.
		\end{enumerate}\\
		\textbf{Postcondición}        & La compra se registra correctamente en \textit{Blockchain} \\
		\textbf{Excepciones}          & No se ha iniciado sesión o no hay fondos \\
		\textbf{Importancia}          & Alta \\
		\bottomrule
	\end{tabularx}
	\caption{CU-8 Compra.}
\end{table}
\apendice{Especificación de diseño}

\section{Introducción}
En este apéndice se va a comentar cómo se ha llevado a cabo el diseño de la aplicación, por ello comentaré cómo se han ido llevando a cabo las distintas partes del diseño.
\section{Diseño de datos}

En primer lugar, el diseño de datos va a ser algo diferente a la forma de hacerlo en \textit{Blockchain}. Por un lado, a bajo nivel el diseño de datos de \textit{Blockchain} se ha explicado en el capítulo de la memoria relativo a conceptos teóricos.


Por ello, en esta sección me voy a centrar en el diseño de datos dentro de los \textit{smart contract} que forman Second Hand Chain que se puede consultar en la figura \ref{fig:anexos/disenioDatos}:

\imagen{anexos/disenioDatos}{Conjunto de datos con los que operan los contratos}

\begin{itemize}
    \item Podemos ver que se compone de un  \textit{struct} que representa un teléfono.
    \item Cada teléfono registrado se va a ir guardando en una lista de teléfonos.
    \item Guardaremos un contador para el número de teléfonos en venta.
    \item En un mapa registraremos a quién pertenece cada teléfono.
    \item Además, guardaremos el número de teléfonos por dueño, este mapa es necesario por el patrón de diseño \textit{memory array building}, explicado en la memoria.
    \item En otro mapa se lleva el recuento de los IMEI ya registrados.
    \item Existe otro mapa para registrar si un teléfono está o no a la venta.
    \item Por otro lado, los dos últimos mapas sirven respectivamente para der permisos a un tercero para transferir un token y para transferir todos los de una cuenta. Estos métodos, son necesarios para poder implementar los métodos encargados de gestionar los permisos de los \textit{NFT} de la interfaz \textit{IERC721}.
    
\end{itemize}

\section{Diseño procedimental}
A continuación se van a mostrar los distintos diagramas de secuencia correspondientes a Second Hand Chain:



%\subsection{Consulta de teléfonos por parte de un usuario sin \textit{wallet}}
\imagen{anexos/cu1seq}{Diagrama de secuencia de la consulta de teléfonos por parte de un usuario sin \textit{wallet}}

%\subsection{Crear un teléfono}
\imagen{anexos/cu3seq}{Diagrama de secuencia de la creación de un teléfono}

%\subsection{Consulta de teléfonos}
\imagen{anexos/cu4seq}{Diagrama de secuencia sobre la consulta de teléfonos}

%\subsection{Puesta a la venta}
\imagen{anexos/cu5seq}{Diagrama de secuencia de puesta a la venta}

%\subsection{Cambio de precio}
\imagen{anexos/cu6seq}{Diagrama de secuencia de cambio de precio de un teléfono a la venta}

%\subsection{Retirada de la venta}
\imagen{anexos/cu7seq}{Diagrama de secuencia de la retirada de la venta de un teléfono}

%\subsection{Compra}
\imagen{anexos/cu8seq}{Diagrama de secuencia de la compra de un teléfono}


\section{Diseño arquitectónico}


En la figura \ref{fig:anexos/despliegue} podemos ver el diagrama de diseño arquitectónico del conjunto de la aplicación:
\begin{itemize}
    \item Los usuarios (tanto de ordenador como de dispositivo móvil)  pueden contar con \textit{Metamask} o directamente con la aplicación.
    \item \textit{Metamask} recibirá solicitudes del \textit{frontend} para firmar transacciones que devolverá a este y serán enviadas. Además un usuario puede iniciar sesión para poder consultar los datos asociados a su cuenta.
    \item \textit{NFT.Storage} recibirá imágenes que almacenará, a cambio devolverá al \textit{frontend} el CID de esa imagen.
    \item \textit{Alchemy}, nuestro proveedor de nodos recibirá peticiones de el \textit{frontend}, a las cuales responderá según tenga esa información. Además ejecutará las transacciones firmadas que reciba del \textit{frontend}.
    \item El \textit{frontend} es la piedra angular de nuestra aplicación, el punto en común entre todos estos servicios y el encargado de su organización.
\end{itemize}

\imagen{anexos/despliegue}{Diagrama de diseño arquitectónico de la aplicación}
\apendice{Documentación técnica de programación}

\section{Introducción}
En primer lugar me gustaría destacar que esta aplicación está completamente desplegada y se puede acceder a ella desde \url{https://second-hand-chain.vercel.app/} donde se pueden verificar todas sus funcionalidades. 

Por ello, para utilizar esta aplicación no es necesario instalar nada (excepto la extensión de \textit{Metamask} en tu navegador).

En este apéndice se va a dar toda la información necesaria al programador para poder configurar la aplicación y sus entornos de desarrollo.

\section{Estructura de directorios}

\begin{itemize}
    \item \textit{documentation}: Documentación.
    \item \textit{Models UML}: Carpeta con los diagramas de la aplicación.
    \item \textit{Truffle}: Contratos Inteligentes.
        \begin{itemize}
            \item \textit{Build/contracts}: ABI\footnote{Interfaz Binaria de Aplicación} de  los contratos.
            \item \textit{contracts}: código de los contratos.
            \item \textit{test}: carpeta creada por defecto para posibles tests .
        \end{itemize}
    \item \textit{frontend/second-hand-chain}: Provecto de \textit{React TypeScript}
            \begin{itemize}
            \item \textit{dist}: Ficheros y videos en el proyecto compilado.
            \item \textit{node\_modules}: node y librerias.
            \item \textit{src}: código.
                \begin{itemize}
                    \item \textit{componentes}: Componentes de \textit{React}.
                    \item \textit{hooks}: \textit{Custom hooks} empleados.
                    \item \textit{services}: Servicios a llamar desde los componentes.
                \end{itemize}
        \end{itemize}
\end{itemize}

\section{Manual del programador}
A continuación se explica el software que es necesario instalar para ejecutar la aplicación:
\subsection{Instalar\textit{Node.js}}
Es necesario instalar la versión v16.14.0 de \textit{Node.js}. Es importante que sea esta versión ya que la mayoría de este proyecto depende de este \textit{software}.

Por otro lado cuenta con la ventaja de ser una versión LTS\footnote{Long term service} para la cual se va a ofrecer soporte durante más tiempo de para otras.

Podemos descargar este software desde su web oficial: 
\url{https://nodejs.dev/en/}.

\subsection{Instalar \textit{Truffle}}

Este entorno de desarrollo para contratos inteligentes nos permitirá crear proyectos de \textit{Truffle} e interactuar con ellos.

Para ello emplearemos el comando: \texttt{npm install -g truffle}.

Además instalaremos  \textit{Ganache} con: \texttt{npm i ganache-cli}

Esta versión de \textit{Ganache} es solamente para consola aunque a mí personalmente me ha parecido más cómoda que la que cuenta con interfaz.

\section{Compilación, instalación y ejecución del proyecto}

\subsection{Obtener claves para  \textit{Alchemy} y \textit{NFT.Storage}}
Para poder ejecutar el proyecto una vez desplegado en \textit{Blockchain} vamos a necesitar las siguientes \textit{API KEYS}:

\begin{itemize}
    \item \textit{Alchemy}: Podemos obtenerla mediante el registro en su web y creando un proyecto, sección en la que nos aparecerá la clave y el \textit{endpoint}. Enlace: https: \url{//docs.alchemy.com/docs/alchemy-quickstart-guide\#1key-create-an-alchemy-key}.
    \item \textit{Nft.Storage}:Deberemos obtener otra clave para este servicio, para lo cual también es necesario el registro. Enlace: \url{https://nft.storage/docs/}.
\end{itemize}

\subsection{Configurar el \textit{.env}}
En los directorios con el \textit{frontend} podemos comprobar que tenemos un fichero llamado \textit{.env.example}, este fichero contiene el nombre exacto de las variables a ser igualadas con las \textit{API KEYS} obtenidas en la sección anterior. Para ello:
\begin{itemize}
    \item Copiar y pegar dicho fichero en el mismo directorio.
    \item La copia renombrarla a \textit{.env}.
    \item Ahora registra las \textit{API KEYS} en ese fichero.
    \item Si tienes un despliegue de los contratos copiar también la dirección del contrato a su variable.
\end{itemize}

\subsection{Instalación de librerías de \textit{Javascript} en los proyectos de \textit{Truffle} y \textit{React}}
Para realizar este paso debemos entrar en el directorio raíz de cada uno de los proyectos y ejecutar: \texttt{npm install}.

Con este comando instalamos todas las librerías de golpe con sus versiones ya configuradas, ya que están definidas en los ficheros de dependencias\footnote{package.json}.

A continuación, para ejecutar el proyecto de \textit{React} deberemos ejecutar en consola: \texttt{npm run dev}.

Tras ello nos aparecerá un enlace en la consola con la dirección local en la que se está ejecutando. No hay que hacer nada más ya que a cada cambio en este proyecto se recompila automáticamente añadiendo solamente el código que ha cambiado gracias a \textit{Vite}.

\subsection{Desplegar el contrato en \textit{Sepolia}}
\begin{itemize}
    \item Descargar el  fichero \textit{truffle-config.js} de uno de los \textit{commits} anteriores a añadirlo a \textit{.gitIgnore}.
    \item Deberemos configurar el fichero \textit{truffle-config.js} siguiendo el formato en la figura \ref{fig:deployment}.
    \item Introducir nuestra clave privada de una \textit{wallet} con suficiente \textit{Ether} donde se indica en la figura \ref{fig:deployment}.
    \item Introducir la \textit{API KEY} de \textit{Alchemy}.
    \item Ejecutar: \texttt{truffle migrate --network sepolia}.
    \item En la salida se esta última ejecución aparece la dirección del contrato, pegarla en su campo correspondiente del fichero \textit{.env}.
\end{itemize}

\imagen{deployment}{Configuración de \textit{truffle-config.js} para despliegue en \textit{Sepolia}}

\section{Pruebas del sistema}
Se ha realizado una serie de pruebas manuales a lo largo del desarrollo de la aplicación para verificar que se cumplen en todo momento los requisitos establecidos.

Se han diseñado estas pruebas para conseguir el mayor porcentaje posible de recubrimiento de código. A continuación se listan cada uno de los casos de prueba testados:

\begin{table}[p]
	\centering
	\begin{tabularx}{\linewidth}{ p{0.21\columnwidth} p{0.71\columnwidth} }
		\toprule
		\textbf{CP-1}    & \textbf{Inicio de Sesión en \textit{Metamask} correcto}\\
		\toprule
		\textbf{Autor}                & Áxel Rubio González \\
		\textbf{Descripción}          & En este caso de prueba el usuario iniciará sesión en  \textit{Metamask} correctamente \\
		\textbf{Pasos}             &
		\begin{enumerate}
			\def\labelenumi{\arabic{enumi}.}
			\tightlist
			\item El usuario accede a la web de Second Hand Chain.
                \item Le aparece una ventana para iniciar sesión en \textit{Metamask}.
                \item Introduce su contraseña.
                \item Puede consultar sus teléfonos en propiedad.
		\end{enumerate}\\
		\textbf{Estado}          & Correcto \\
		\bottomrule
	\end{tabularx}
	\caption{CP-1 Inicio de Sesión en \textit{Metamask} correcto.}
\end{table}


\begin{table}[p]
	\centering
	\begin{tabularx}{\linewidth}{ p{0.21\columnwidth} p{0.71\columnwidth} }
		\toprule
		\textbf{CP-2}    & \textbf{Inicio de Sesión en \textit{Metamask} pero cierra la ventana antes}\\
		\toprule
		\textbf{Autor}                & Áxel Rubio González \\
		\textbf{Descripción}          & En este caso de prueba el usuario accederá a la web, cerrará la ventana de \textit{Metamask} y accederá desde la extensión del navegador. \\
		\textbf{Pasos}             &
		\begin{enumerate}
			\def\labelenumi{\arabic{enumi}.}
			\tightlist
			\item El usuario accede a la web de Second Hand Chain.
                \item Le aparece una ventana para iniciar sesión en \textit{Metamask}.
                \item Cierra la Ventana.
                \item Le aparece un mensaje diciéndole que acceda desde la extensión del navegador en el botón derecho de la \textit{NavBar}.
                \item Accede a la extensión del navegador de \textit{Metamask}.
                \item Introduce su contraseña.
                \item Su cuenta aparece en el botón.
                \item Puede consultar sus teléfonos en propiedad.
		\end{enumerate}\\
		\textbf{Estado}          & Correcto \\
		\bottomrule
	\end{tabularx}
	\caption{CP-2 Inicio de Sesión en \textit{Metamask} pero cierra la ventana antes.}
\end{table}


\begin{table}[p]
	\centering
	\begin{tabularx}{\linewidth}{ p{0.21\columnwidth} p{0.71\columnwidth} }
		\toprule
		\textbf{CP-3}    & \textbf{Usuarios sin \textit{Metamask}}\\
		\toprule
		\textbf{Autor}                & Áxel Rubio González \\
		\textbf{Descripción}          & En este caso de prueba el usuario accederá a la web, sin la extensión instalada, podrá consultar datos pero no operar. \\
		\textbf{Pasos}             &
		\begin{enumerate}
			\def\labelenumi{\arabic{enumi}.}
			\tightlist
			\item El usuario accede a la web de Second Hand Chain.
                \item Le aparece un mensaje diciéndole que instale \textit{Metamask}.
                \item Accede a consultar los teléfonos en venta.
                \item Selecciona uno.
                \item Le aparece un mensaje en el botón de compra diciendole que inicie sesión en \textit{Metamask}.
		\end{enumerate}\\
		\textbf{Estado}          & Correcto \\
		\bottomrule
	\end{tabularx}
	\caption{CP-3 Usuarios sin \textit{Metamask}.}
\end{table}


\begin{table}[p]
	\centering
	\begin{tabularx}{\linewidth}{ p{0.21\columnwidth} p{0.71\columnwidth} }
		\toprule
		\textbf{CP-4}    & \textbf{Creación correcta de teléfono}\\
		\toprule
		\textbf{Autor}                & Áxel Rubio González \\
		\textbf{Descripción}          & En este caso de prueba el usuario, con la cuenta de \textit{Metamask} ya configurada  registrará correctamente un teléfono. \\
		\textbf{Pasos}             &
		\begin{enumerate}
			\def\labelenumi{\arabic{enumi}.}
			\tightlist
			\item El usuario accede a la web de Second Hand Chain.
                \item Hace clic en crear teléfono.
                \item Introduce todos los datos correctamente, incluyendo la imagen que selecciona de su dispositivo.
                \item Hace clic en crear.
                \item Firma correctamente la transacción en \textit{Metamask}.
                \item Le aparece un mensaje con un enlace a Etherscan.
                \item Accede satisfactoriamente a \textit{Etherscan} para seguir el estado de su transacción.
		\end{enumerate}\\
		\textbf{Estado}          & Correcto \\
		\bottomrule
	\end{tabularx}
	\caption{CP-4 Creación correcta de teléfono.}
\end{table}

\begin{table}[p]
	\centering
	\begin{tabularx}{\linewidth}{ p{0.21\columnwidth} p{0.71\columnwidth} }
		\toprule
		\textbf{CP-5}    & \textbf{Creación incorrecta de teléfono por datos incompletos}\\
		\toprule
		\textbf{Autor}                & Áxel Rubio González \\
		\textbf{Descripción}          & En este caso de prueba el usuario, con la cuenta de \textit{Metamask} ya configurada  intentará registrar un teléfono sin haber introducido todos los datos necesarios. \\
		\textbf{Pasos}             &
		\begin{enumerate}
			\def\labelenumi{\arabic{enumi}.}
			\tightlist
			\item El usuario accede a la web de Second Hand Chain.
                \item Hace clic en crear teléfono.
                \item Introduce todos los datos excepto uno que no completa.
                \item Hace clic en crear.
                \item Le aparece un mensaje avisando de que tiene datos incompletos.
		\end{enumerate}\\
		\textbf{Estado}          & Correcto \\
		\bottomrule
	\end{tabularx}
	\caption{CP-5 Creación incorrecta de teléfono por datos incompletos.}
\end{table}


\begin{table}[p]
	\centering
	\begin{tabularx}{\linewidth}{ p{0.21\columnwidth} p{0.71\columnwidth} }
		\toprule
		\textbf{CP-6}    & \textbf{Creación incorrecta de teléfono por IMEI ya registrado}\\
		\toprule
		\textbf{Autor}                & Áxel Rubio González \\
		\textbf{Descripción}          & En este caso de prueba el usuario, con la cuenta de \textit{Metamask} ya configurada,  intentará registrar un teléfono previamente registrado. \\
		\textbf{Pasos}             &
		\begin{enumerate}
			\def\labelenumi{\arabic{enumi}.}
			\tightlist
			\item El usuario accede a la web de Second Hand Chain.
                \item Hace clic en crear teléfono.
                \item Introduce todos los datos, pero el IMEI pertenece a otro teléfono ya registrado.
                \item Hace clic en crear.
                \item Le aparece un mensaje avisando de que ese IMEI ya existe en el contrato.
		\end{enumerate}\\
		\textbf{Estado}          & Correcto \\
		\bottomrule
	\end{tabularx}
	\caption{CP-6 Creación incorrecta de teléfono por IMEI ya registrado.}
\end{table}

\begin{table}[p]
	\centering
	\begin{tabularx}{\linewidth}{ p{0.21\columnwidth} p{0.71\columnwidth} }
		\toprule
		\textbf{CP-7}    & \textbf{Puesta a la venta correcta de teléfono}\\
		\toprule
		\textbf{Autor}                & Áxel Rubio González \\
		\textbf{Descripción}          & En este caso de prueba, con la cuenta de \textit{Metamask} ya configurada, el usuario pondrá a la venta un teléfono de su propiedad correctamente. \\
		\textbf{Pasos}             &
		\begin{enumerate}
			\def\labelenumi{\arabic{enumi}.}
			\tightlist
			\item El usuario accede a la web de Second Hand Chain.
                \item Hace clic en teléfonos en propiedad.
                \item Selecciona uno.
                \item Introduce el precio.
                \item Hace clic en poner a la venta.
                \item Firma correctamente la transacción en \textit{Metamask}.
                \item Le aparece un mensaje con un enlace a \textit{Etherscan}.
                \item Accede satisfactoriamente a \textit{Etherscan} para seguir el estado de su transacción.
                \item Espera a que la transacción se confirme.
                \item Accede a teléfonos en venta y puede ver el teléfono a la venta.
		\end{enumerate}\\
		\textbf{Estado}          & Correcto \\
		\bottomrule
	\end{tabularx}
	\caption{CP-7 Puesta a la venta correcta de teléfono.}
\end{table}

\begin{table}[p]
	\centering
	\begin{tabularx}{\linewidth}{ p{0.21\columnwidth} p{0.71\columnwidth} }
		\toprule
		\textbf{CP-8}    & \textbf{Retirada de la venta correcta de teléfono}\\
		\toprule
		\textbf{Autor}                & Áxel Rubio González \\
		\textbf{Descripción}          & En este caso de prueba el usuario, con la cuenta de \textit{Metamask} ya configurada, retirará un teléfono de la venta. \\
		\textbf{Pasos}             &
		\begin{enumerate}
			\def\labelenumi{\arabic{enumi}.}
			\tightlist
			\item El usuario accede a la web de Second Hand Chain.
                \item Hace clic en teléfonos en propiedad.
                \item Selecciona uno que este a la venta.
                \item Hace clic en retirar de la venta.
                \item Firma correctamente la transacción en \textit{Metamask}.
                \item Le aparece un mensaje con un enlace a \textit{Etherscan}.
                \item Accede satisfactoriamente a \textit{Etherscan} para seguir el estado de su transacción.
                \item Espera a que la transacción se confirme.
                \item Accede a teléfonos en venta y puede ver el teléfono no se puede encontrar.
                \item Accede a teléfonos en propiedad
                \item Selecciona ese teléfono.
                \item No aparece la opción de retirar de la venta.
		\end{enumerate}\\
		\textbf{Estado}          & Correcto \\
		\bottomrule
	\end{tabularx}
	\caption{CP-8 Retirada de la venta correcta de teléfono.}
\end{table}

\begin{table}[p]
	\centering
	\begin{tabularx}{\linewidth}{ p{0.21\columnwidth} p{0.71\columnwidth} }
		\toprule
		\textbf{CP-9}    & \textbf{Cambio de precio de teléfono a la venta}\\
		\toprule
		\textbf{Autor}                & Áxel Rubio González \\
		\textbf{Descripción}          & En este caso de prueba el usuario, con la cuenta de \textit{Metamask} ya configurada, cambiará el precio de un teléfono la venta. \\
		\textbf{Pasos}             &
		\begin{enumerate}
			\def\labelenumi{\arabic{enumi}.}
			\tightlist
			\item El usuario accede a la web de Second Hand Chain.
                \item Hace clic en teléfonos en propiedad.
                \item Selecciona uno que esté a la venta.
                \item Escribe un nuevo precio.
                \item Hace clic en cambiar precio.
                \item Firma correctamente la transacción en \textit{Metamask}.
                \item Le aparece un mensaje con un enlace a \textit{Etherscan}.
                \item Accede satisfactoriamente a \textit{Etherscan} para seguir el estado de su transacción.
                \item Espera a que la transacción se confirme.
                \item Accede a teléfonos en venta y puede ver el teléfono con su nuevo precio.
		\end{enumerate}\\
		\textbf{Estado}          & Correcto \\
		\bottomrule
	\end{tabularx}
	\caption{CP-9 Cambio de precio de teléfono.}
\end{table}

\begin{table}[p]
	\centering
	\begin{tabularx}{\linewidth}{ p{0.21\columnwidth} p{0.71\columnwidth} }
		\toprule
		\textbf{CP-10}    & \textbf{Compra de Teléfono}\\
		\toprule
		\textbf{Autor}                & Áxel Rubio González \\
		\textbf{Descripción}          & En este caso de prueba el usuario, con la cuenta de \textit{Metamask} ya configurada, comprará un teléfono. \\
		\textbf{Pasos}             &
		\begin{enumerate}
			\def\labelenumi{\arabic{enumi}.}
			\tightlist
			\item El usuario accede a la web de Second Hand Chain.
                \item Hace clic en teléfonos a la venta.
                \item Selecciona uno que esté a la venta.
                \item Hace clic en comprar.
                \item Firma correctamente la transacción en \textit{Metamask}.
                \item Le aparece un mensaje con un enlace a \textit{Etherscan}.
                \item Accede satisfactoriamente a \textit{Etherscan} para seguir el estado de su transacción.
                \item Espera a que la transacción se confirme.
                \item Accede a teléfonos en propiedad y puede ver el teléfono recién adquirido.
		\end{enumerate}\\
		\textbf{Estado}          & Correcto \\
		\bottomrule
	\end{tabularx}
	\caption{CP-10 Compra de Teléfono.}
\end{table}


\apendice{Documentación de usuario}

\section{Introducción}
En este anexo se explica todo lo que necesita el usuario para utilizar la aplicación, y se explica el funcionamiento de la misma, aunque también cuente como soporte con los vídeos de la aplicación.

\section{Requisitos de usuarios}
A continuación, se exponen los requisitos necesarios para poder utilizar la aplicación:
\begin{itemize}
    \item Tener instalada la extensión de navegador \textit{Metamask}.
    \item Es preferible usar \textit{Brave} como navegador aunque al estar basado en \textit{Chromium} no hay problemas en emplear cualquier otro que emplee esa base de código abierto.
    \item Para usuarios de dispositivos móviles tener instalada la aplicación \textit{Metamask}.
\end{itemize}

\section{Instalación}
Los usuarios deberán añadir la extensión \textit{Metamask} a su navegador, en la web desplegada hay un vídeo con el proceso de instalación, también se encuentra en el repositorio y en los USB entregados, en los ficheros públicos de la aplicación web.

Por otro lado, para poder realizar transacciones los usuarios necesitan \textit{Ether} de la red \textit{Sepolia}, pueden obtenerlos registrándose en \textit{Alchemy} y recurriendo a su \textit{Faucet} (como se muestra en el segundo vídeo en la página web) o los pueden pedir al autor por correo aportando la dirección de su \textit{wallet}.

\section{Manual del usuario}
Finalmente en esta sección vamos a detallar el funcionamiento de la aplicación Second Hand Chain para que los usuarios puedan comprender su funcionamiento detalladamente:

\imagen{anexos/vistaInicial}{Vista inicial de la aplicación.}

En la figura \ref{fig:anexos/vistaInicial} podemos destacar los siguientes apartados:
\begin{enumerate}
    \item Botón que nos lleva a la vista actual de presentación de la aplicación.
    \item Botón que abre el modal para registrar teléfonos.
    \item Botón para acceder a los teléfonos que posee el usuario.
    \item Botón para acceder a los teléfonos en venta.
    \item Botón que indica el estado de conexión con \textit{Metamask}.
    \item Ventana para iniciar sesión en \textit{Metamask}, aparece si acabas de abrir tu navegador y tienes \textit{Metamask} instalado, permite iniciar sesión, proceso necesario para poder firmar transacciones.
    \item Descripción de la aplicación.
    \item Vídeo para mostrar como añadir \textit{Metamask}.
    \item Vídeo para mostrar como añadir \textit{Ether} a \textit{Metamask} gracias al \textit{faucet} de \textit{Alchemy} y demostración del funcionamiento de la aplicación.
    \item Botón para acceder al repositorio de \textit{GitHub}.
    \item Listado de tecnologías empleadas en el proyecto.
\end{enumerate}

\imagen{anexos/vistatarjeta}{Vista de tarjetas de teléfonos.}

En la figura \ref{fig:anexos/vistatarjeta} podemos ver un par de tarjetas de teléfonos:
\begin{enumerate}
    \item Modelo y marca del teléfono.
    \item Características del teléfono.
    \item Botón de compra, permite acceder a los detalles, en el caso de los teléfonos en propiedad aparecerá simplemente la palabra detalles.
    \item Características del teléfono.
    
\end{enumerate}

\imagen{anexos/vistadetalles}{Vista de detalles de teléfonos.}

En la figura \ref{fig:anexos/vistadetalles} podemos ver un modal con los detalles del teléfono:
\begin{enumerate}
    \item Modelo y marca del teléfono.
    \item Características del teléfono.
    \item Historial de acciones de compraventa, siendo la de índice cero la del registro del teléfono.
    \item Botón de compra, puede aparecer deshabilitado si no iniciamos sesión, habilitado para poderlo comprar o no aparecer, en el caso de acceder desde teléfonos en propiedad.
\end{enumerate}

\imagen{anexos/vistadetalles2}{Vista de detalles de teléfonos en propiedad, parte para configurar el estado del teléfono.}

En la figura \ref{fig:anexos/vistadetalles2} podemos ver un segmento del modal con las opciones a configurar por un propietario para poner a la venta el teléfono:
\begin{enumerate}
    \item Si el teléfono está a la venta o no y su precio, en caso de estarlo.
    \item Campo para introducir el precio por el teléfono.
    \item Botón de cambiar de precio si está a la venta o de poner en venta si no lo estaba.
    \item Botón de retirar de la venta, solo aparece si se vende el teléfono.
\end{enumerate}

\imagen{anexos/vistatransaccion}{Vista de transacción de compra de teléfono.}

En la figura \ref{fig:anexos/vistatransaccion} podemos ver el proceso para comprar un teléfono:
\begin{enumerate}
    \item Se hace clic en el botón de compra por el cual nos aparece un modal de \textit{Metamask} solicitando que firmemos la transacción.
    \item Podemos ver la cuenta con la que vamos a firmar la transacción.
    \item Estamos en la \textit{testnet Sepolia}.
    \item El método que vamos a ejecutar, en este caso es detectado al pertenecer a \textit{ERC721}.
    \item Valor de \textit{Ether} que vamos a aportar.
    \item \textit{Gas fees}.
    \item \textit{Valor total de \textit{Ether} que se va a emplear.}
\end{enumerate}
Finalmente confirmaríamos la transacción.

\imagen{anexos/viewlink}{Vista de transacción de compra de teléfono.}

En la figura \ref{fig:anexos/viewlink} podemos ver el mensaje con un enlace para abrir la transacción realizada en \textit{Etherscan}.

\imagen{anexos/viewetherscan}{Vista de la transacción de compra realizada en \textit{Etherscan}.}

La figura \ref{fig:anexos/viewetherscan} nos muestra los detalles de la transacción previamente realizada, en \textit{Etherscan}, vamos a ver sus detalles:
\begin{enumerate}
    \item Nos recuerda que la transacción pertenece a la \textit{testnet Sepolia}.
    \item El \textit{hash} de la transacción.
    \item Estamos en la \textit{testnet Sepolia}.
    \item Momento en el que se minó el bloque al que pertenece nuestra transacción.
    \item Dirección origen, la del comprador.
    \item Dirección con la que ha interactuado la anterior, la de nuestro contrato inteligente.
    \item Nos detecta que es un token en cuyo contrato inteligente se implementa la interfaz \textit{IERC721}
    \item Identificador del token transferido.
    \item Descripción del contrato.
    \item Antiguo propietario del \textit{NFT}.
    \item Nuevo propietario del \textit{NFT}.
    \item Valor en \textit{Ether} transferido.
\end{enumerate}


\bibliographystyle{plain}
\bibliography{bibliografiaAnexos}

\end{document}
